\documentclass[10pt]{article}
\usepackage[a4paper,margin=1.4cm]{geometry}
\usepackage[T1]{fontenc}
\usepackage{lmodern}
\usepackage{booktabs}
\usepackage{multirow}
\usepackage[table]{xcolor}
\usepackage{caption}
\captionsetup{font=small,labelfont=bf,skip=4pt}
\setlength{\parskip}{3pt}
\setlength{\parindent}{0pt}
\pagestyle{empty}

\definecolor{ours}{gray}{0.92}
\newcommand{\gr}[1]{\textcolor{gray}{\textit{#1}}}

\begin{document}

\begin{center}
{\large\bfseries V-JEPA 2.1 on IntPhys 2 (Main split)}
\end{center}

\begin{table}[h]
\centering
\caption{\textbf{Accuracy values showing best model performance for each difficulty level.}
Rows other than V-JEPA 2.1 are reproduced from IntPhys 2 (Bordes et al., 2025), Table 2: each model's best run per
column over a hyperparameter grid, with surprise averaged over time (the paper's default, App.~D.1). For
V-JEPA 2.1 (shaded, ours), each column is the best of the 6 context lengths; the max-surprise row is given for
reference only. Bold: best per model type, ours included.}
\small
\begin{tabular}{l|c|ccc|c||c|c}
\toprule
Model & Type & Easy & Medium & Hard & Overall & Held Out & \gr{IntPhys} \\
\midrule
GPT4-o & MLLMs & 57.69 & 54.75 & 54.17 & 53.75 & 53.19 & \gr{-} \\
Qwen-VL 2.5 & MLLMs & 50.96 & 53.25 & 51.49 & 52.27 & 49.12 & \gr{-} \\
Gemini-1.5 Pro & MLLMs & 58.65 & 53.0 & 52.67 & 52.27 & 52.10 & \gr{55.81} \\
Gemini-2.5 Flash & MLLMs & \textbf{64.42} & \textbf{56.75} & \textbf{54.46} & \textbf{55.63} & \textbf{56.10} & \gr{56.39} \\
\midrule
VideoMAEv2-g & Predictive & 46.00 & \textbf{58.50} & 52.73 & 53.75 & 53.49 & \gr{59.40} \\
Cosmos-4B & Predictive & 46.00 & 52.00 & 48.05 & 49.41 & 48.84 & \gr{85.42} \\
V-JEPA-h + RoPE & Predictive & 52.00 & 53.00 & 57.42 & 53.75 & 54.65 & \gr{98.30} \\
V-JEPA 2-h & Predictive & 54.00 & \textbf{58.50} & \textbf{59.38} & \textbf{57.51} & \textbf{56.40} & \gr{87.22} \\
\rowcolor{ours}
V-JEPA 2.1-g (avg surprise) & Predictive & \textbf{63.46} & 55.50 & 55.95 & 55.14 & -- & \gr{-} \\
\rowcolor{ours}
\quad\textit{max surprise, ref.} & Predictive & 50.00 & 52.00 & 54.76 & 51.78 & -- & \gr{-} \\
\midrule
Human & - & 96.17 & 97.8 & 95.5 & 96.44 & 92.44 & \gr{-} \\
\bottomrule
\end{tabular}
\end{table}

\begin{table}[h]
\centering
\caption{\textbf{Accuracy across property and camera type.} For each model, the accuracy of each subset is the best one
across a set of hyperparameters (for V-JEPA 2.1: the 6 context lengths, avg surprise; max surprise for reference). The smaller size of each subset contributes to
volatility in performance (52 / 68 pairs per Fixed / Moving cell, 94 for Solidity Moving).}
\small
\setlength{\tabcolsep}{4pt}
\begin{tabular}{l|cc|cc|cc|cc}
\toprule
& \multicolumn{2}{c|}{Permanence} & \multicolumn{2}{c|}{Immutability} & \multicolumn{2}{c|}{Continuity} & \multicolumn{2}{c}{Solidity} \\
Model & Fixed & Moving & Fixed & Moving & Fixed & Moving & Fixed & Moving \\
\midrule
GPT4-o & 59.62 & \textbf{58.82} & 58.65 & 59.56 & \textbf{54.81} & \textbf{57.35} & \textbf{56.73} & 55.32 \\
Qwen-VL 2.5 & 53.85 & 54.41 & 56.73 & 53.68 & 52.88 & 54.41 & 50.96 & 51.06 \\
Gemini-1.5 Pro & 55.77 & 55.88 & 56.73 & 56.73 & 54.80 & 54.80 & 56.73 & 56.73 \\
Gemini-2.5 Flash & \textbf{64.42} & \textbf{58.82} & \textbf{59.62} & \textbf{63.97} & \textbf{54.81} & 55.15 & 55.77 & \textbf{56.38} \\
\midrule
VideoMAEv2-g & 63.46 & 50.00 & 54.81 & 53.69 & \textbf{65.38} & 54.41 & 48.08 & \textbf{59.57} \\
Cosmos-4B & 51.92 & 41.18 & 50.96 & 48.32 & 53.85 & 50.00 & 48.08 & 55.32 \\
V-JEPA-h + RoPE & 63.46 & \textbf{67.65} & 51.92 & 56.38 & 50.00 & 57.35 & 50.00 & 52.13 \\
V-JEPA 2-h & 59.62 & 57.35 & \textbf{55.77} & \textbf{58.72} & 57.69 & \textbf{75.00} & 46.15 & 58.51 \\
\rowcolor{ours}
V-JEPA 2.1-g (avg surprise) & \textbf{67.31} & 54.41 & \textbf{55.77} & 54.41 & 59.62 & 60.29 & \textbf{53.85} & 55.32 \\
\rowcolor{ours}
\quad\textit{max surprise, ref.} & 44.23 & 55.88 & 51.92 & 50.00 & 59.62 & 52.94 & 51.92 & 53.19 \\
\midrule
Human & 100.0 & 99.26 & 97.11 & 90.44 & 99.04 & 94.44 & 96.15 & 95.21 \\
\bottomrule
\end{tabular}
\end{table}

\vspace{-2pt}
\small
\textbf{Experimental setup.}
\emph{Model.} Official V-JEPA 2.1 ViT-g/16 release at 384\,px from facebookresearch/\allowbreak vjepa2
(checkpoint \texttt{vjepa2\_1\_vitg\_384.pt}). The model code is vendored unchanged from that repository, and the context encoder,
target (EMA) encoder (1.01\,B parameters each) and predictor (55\,M) load with strict key matching. V-JEPA~2.1
predicts the concatenation of 4 hierarchical encoder levels; we use only the last one (layer 39), layer-normed
as in its training targets, and the predictor's mask token 0 (the one used for video during training).

\emph{Surprise.} Following Garrido et al.\ (2025) and the setup reported as best for V-JEPA~2 in IntPhys~2:
videos are sampled every 10th frame (6\,fps), split into 48-frame windows with stride 2, and for each window the
model predicts the features of all frames after a context of $C\in\{12,18,24,30,36,42\}$ frames, as in the
authors' released V-JEPA~2 config (\texttt{vjepa\_2.yaml}; the paper lists $\{4,\dots,14\}$). The prediction error is
the L1 distance between predicted and target features (bf16). With \texttt{max\_context\_mode} (the paper's
growing-context protocol), the first window is also predicted from every shorter context $2,4,\dots,C-2$, so that
early events are not missed. A video's surprise is the average of these errors over time (App.~D.1).

\emph{Hyperparameters.} For V-JEPA and V-JEPA~2 the paper reports, per column, the maximum over window size
$\{16,32,48\}$ at a fixed 6\,fps, prediction length (fixed, or to the end of the window) and context length
(App.~D.3); V-JEPA~2's best was window 48, prediction to the end of the window, 6\,fps (Table~8). We evaluate only
that setting and maximise over the 6 context lengths, a subset of their search.

\emph{Metric.} Relative accuracy on matched pairs: in each scene quadruplet, the impossible video
\texttt{k\_Impossible} must be more surprising than its matched \texttt{k\_Possible} ($k=1,2$), giving 506 pairs
on Main (1,012 videos), with 95\% confidence intervals of about $\pm$4.4 points overall and $\pm$14 on Easy.
Easy/Medium/Hard follow the environment lists of the authors' scoring notebook (52/200/168 pairs; the other 86
pairs count only towards Overall). The paper's Easy and Hard values are not multiples of $1/52$ and $1/168$, so
its subsets differ slightly. Held Out was not evaluated.

\emph{Observation.} With the paper's average surprise, V-JEPA~2.1-g reaches 55.14\% overall, below
V-JEPA~2-h (57.51\%). Max surprise stays at chance (51.78\%): for 81--97\% of videos the maximum comes from the
short-context predictions on the first window, so context length barely matters.

\emph{Compute.} 6 single-GPU nodes (RTX 5090, 3090, 5070\,Ti, 5060\,Ti, 4060\,Ti, plus one more 16\,GB card),
$\approx$21\,h wall-clock, 1.3--7.3\,min per video depending on the GPU; videos split evenly across ranks, with no
inter-node communication.

\end{document}
